\subsection{Examples}

In this number, we assume X is pure of finite dimension n. We also assume it is endowed with a manifold structure of class \(C^{r+1}\) compatible with the given structure of class \(C^{r}\) .

14.4.1 ("Infinitesimal transformations"). Let \(\tau\) be a finite-dimensional vector functor for isomorphisms, of class \(C^r\) ; we assume that, for every finite-dimensional vector space V, \(\tau(V)\) is of finite dimension. We denote by \(E_t\) the vector bundle \(\tau(T(X))\) ; it is of class \(C^r\) .

Let \(\xi\) be a vector field of class \(C^r\) on X . There exists a differential operator D of type \(E_{\tau} \to E_{\tau}\) and of order \(\leqslant1\) such that

\[
\mathbf {D} _ {\mathrm{U}} (s) = \theta_ {\xi}. s
\]

for every open set U of X and every \(s \in \mathcal{S}_{\mathbb{E}}^{r}(\mathrm{U})\) , cf. 8.4.3. Such an operator D is unique; we denote it by \((\theta_{\xi})_{\tau}\) or simply \(\theta_{\xi}\) . Its symbol is the section \(\xi \otimes \mathrm{Id}_{\mathbb{E}_{\tau}}\) of the vector bundle

\[
\mathrm{S} ^ {1} (\mathrm{E} _ {\tau}, \mathrm{E} _ {\tau}) = \mathrm{T} (\mathrm{X}) \otimes \mathscr {L} (\mathrm{E} _ {\tau}; \mathrm{E} _ {\tau}).
\]

If V is a finite-dimensional K-vector space, set

\[
\tilde {\tau} (\mathbf {V}) = \mathscr {L} (\tau (\mathbf {V}); \mathrm{Alt} ^ {n} (\mathbf {V}; \mathbf {K})) = \tau (\mathbf {V}) ^ {*} \otimes \bigwedge^ {n} \mathbf {V} ^ {*},
\]

and if \(u\colon V_{1}\to V_{2}\) is an isomorphism, define \(\tilde{\tau}(u)\colon\tilde{\tau}(V_{1})\to\tilde{\tau}(V_{2})\) by transport of structure. One thus obtains a vector functor in finite dimension \(\tilde{\tau}\) . The bundle \(E_{\tilde{\tau}}=\tilde{\tau}(T(X))\) is identified in an obvious manner with the bundle \((E_{\tau})^{\sim}\) ; the transpose of \((\theta_{\xi})_{\tau}\) is \(-( \theta_{\xi})_{\tilde{\tau}}\) and formula (1) of no. 14.3.4 takes the form

\[
\langle \theta_ {\xi}. u, v \rangle + \langle u, \theta_ {\xi}. v \rangle = d (i (\xi) (u. v)).
\]

14.4.2 (“Exterior differential”). For every integer \(p \geqslant 0\) one sets

\[
\Omega^ {p} = \operatorname{Alt} ^ {p} (\mathrm{T} (\mathrm{X}); \mathrm{K} _ {\mathrm{X}}).
\]

There exists a differential operator D of type \(\Omega^{p} \to \Omega^{p+1}\) and of order \(\leqslant 1\) such that

\[
\mathrm{D} _ {\mathrm{U}} (\omega) = d \omega
\]

for every open set U of X and every \(\omega\in\mathcal{S}_{\Omega^{p}}^{r}(U)\) . Such an operator D is unique; it is denoted \(d\) (or \(d_{p}\) if one wishes to specify the integer \(p\) ).

Its symbol is the element of \(S^1(\Omega^p, \Omega^{p+1})\) obtained in the following manner: one first has

\[
\begin{array}{r l} \mathrm{S} ^ {1} (\Omega^ {p}, \Omega^ {p + 1}) & = \mathrm{T} (\mathrm{X}) \otimes \mathcal {L} (\Omega^ {p}; \Omega^ {p + 1}) \\ & = \mathcal {L} (\Omega^ {1}; \mathcal {L} (\Omega^ {p}; \Omega^ {p + 1})) = \mathcal {L} (\Omega^ {1}, \Omega^ {p}; \Omega^ {p + 1}), \end{array}
\]

where \(\mathcal{L}(\Omega^{1},\Omega^{p};\Omega^{p+1})\) denotes the bundle of bilinear maps from \(\Omega^{1}\times\Omega^{p}\) to \(\Omega^{p+1}\) . Now the exterior product \((\alpha,\beta)\mapsto\alpha\wedge\beta\) defines a canonical section of this latter bundle; this section is the symbol of \(d_{p}\) .

To determine the transpose of \(d_{p}\) , one first notes that the exterior product defines a pairing \(\Omega^{p} \times \Omega^{n-p} \to \Omega^{n} = \Omega\) which allows one to identify \((\Omega^{p})^{\sim}\) with \(\Omega^{n-p}\) . One identifies in the same manner \((\Omega^{p+1})^{\sim}\) with \(\Omega^{n-p-1}\) . With these conventions, the transpose of \(d_{p}\) is \((-1)^{p+1}d_{n-p-1}\) and the corresponding Green operator, of type \((\Omega^{p}, \Omega^{n-p-1}) \to \Omega_{1} = \Omega^{n-1}\) , is simply the exterior product (on each fiber).

14.4.3 (“Laplacian”). One takes K = R, r = ∞, X = R \(^{n}\) , E = F = K \(_{X}\) ; one denotes \(\xi^{i}\) (1 ≤ i ≤ n) the coordinate functions on X; one sets D \(_{i}\) = ∂/∂ξ \(^{i}\) , and L = ∑ \(_{i=1}^{n}\) D \(_{i}^{2}\) . The operator L is a scalar differential operator of order ≤ 2. If one identifies Ω with K \(_{X}\) by means of the frame ω = dξ \(^{1}\) ∧ ⋯ ∧ dξ \(^{n}\) , one has \(\tilde{E}\) = \(\tilde{F}\) = K \(_{X}\) (cf. 14.3.3, example) and \(^{t}\) L = L.

For \(1 \leqslant i \leqslant n\) , set \(\omega_{i} = (-1)^{i-1} d\xi^{1} \wedge \cdots \wedge d\xi^{i-1} \wedge d\xi^{i+1} \wedge \cdots \wedge d\xi^{n}\) . The \(\omega_{i}\) form a frame of \(\Omega_{1} = \Omega^{n-1}\) . If \(f \in \mathcal{C}^{1}(X)\) , denote by \(\text{grad}(f)\) the section \(\sum_{i=1}^{n} D_{i}(f)\omega_{i}\) of \(\Omega_{1}\) . There exists a Green operator G for L such that one has

\(\mathbf{G}(u,v) = v.\operatorname{grad}(u) - u.\operatorname{grad}(v)\) for every open set U of X and \(u,v \in \mathcal{C}^{1}(\mathbf{U})\) .

In particular, let A be a compact piece of \(\mathbf{R}^{n}\) . Orient \(\mathbf{R}^{n}\) , endow A with the orientation induced by that of \(\mathbf{R}^{n}\) and \(\partial\mathbf{A}\) with the corresponding orientation (11.2.1). One has

\[
\int_ {\mathbf {A}} (\mathbf {L} u). v \omega - \int_ {\mathbf {A}} u. (\mathbf {L} v) \omega = \int_ {\partial \mathbf {A}} (v. \operatorname{grad} (u) - u. \operatorname{grad} (v))
\]

for u, v in \(\mathcal{C}^{1}(X)\) .

